\section{A fonte de surto (1,2/50\,\texorpdfstring{$\mu$}{u}s)}
\label{sec:fonte}

A excitação é o impulso atmosférico padrão, modelado por uma dupla exponencial:

\begin{equation}
\Vsrc = V_{\mathrm{pico}}\,K\left(e^{-a t} - e^{-b t}\right), \qquad t \ge 0,
\end{equation}

em que \(a = \ln(2)/T_2\) controla a cauda e \(b = 3/T_1\) controla a frente. Para
\(T_1 = \SI{1.2}{\micro\second}\) (tempo de frente) e \(T_2 = \SI{50}{\micro\second}\) (tempo de cauda),
tem-se \(a \approx \SI{1.3863e4}{s^{-1}}\) e \(b \approx \SI{2.5e6}{s^{-1}}\); a constante \(K\)
normaliza o pico para \(V_{\mathrm{pico}} \approx \SI{1}{kV}\). É a forma clássica de ensaio de impulso
de alta tensão \cite{kuffel2000,iec60060}.

O que torna esse surto crítico para um enrolamento não é a amplitude, e sim a \emph{frente rápida}: os
\SI{1.2}{\micro\second} de subida injetam um espectro de alta frequência que enxerga o enrolamento como
uma rede capacitiva. É essa frente que produz a distribuição inicial não uniforme analisada nas seções
seguintes; a cauda de \SI{50}{\micro\second} apenas sustenta a tensão enquanto a redistribuição
interna ocorre.
